% sections/theory.tex -- orthogonality reformulation, hierarchy, swap formula, inertia.
% Exact statements checked by Analysis/revision/theory/verify_theory.py.

\section{Orthogonality conditions on the height function}
\label{sec:theory}

Throughout, $n\ge 3$, $N=n^{2}$, and a cell of the $n\times n$ grid is a pair
$c=(i,j)$ with $0\le i,j\le n-1$; an arrangement is a bijection $S=(s_{ij})$ from the cells onto
$\{1,\dots,N\}$ (\Cref{def:magic_square}).  We use the centred
coordinates of \Cref{def:gem_coords},
\[
x_{ij}=j-\tfrac{n-1}{2},\qquad y_{ij}=\tfrac{n-1}{2}-i,\qquad
z_{ij}=s_{ij}-\tfrac{N+1}{2},
\]
and refer to $z=z(S)$ as the \emph{height function} of $S$.
A \emph{line} is one of the $n$ rows, the $n$ columns, the main diagonal
$\{i=j\}$ or the anti-diagonal $\{i+j=n-1\}$; there are $2n+2$ lines, each
containing exactly $n$ cells.  With $M(n)=n(N+1)/2$ the magic constant, the
\emph{deviation} of a line $L$ is
\begin{equation}
\delta_L(S)\;=\;\sum_{c\in L}s_c-M(n)\;=\;\sum_{c\in L}z_c ,
\label{eq:line_dev}
\end{equation}
the second equality because $|L|=n$ and $M(n)=n\cdot\frac{N+1}{2}$.
$S$ is \emph{semi-magic} if $\delta_L=0$ for the $2n$ rows and columns, and
\emph{magic} if $\delta_L=0$ for all $2n+2$ lines.  We write $\rho_k=\delta_{\text{row }k}$,
$\gamma_\ell=\delta_{\text{col }\ell}$, $\delta_{\mathrm{m}}$ and $\delta_{\mathrm{a}}$ for the
row, column, main-diagonal and anti-diagonal deviations.

\subsection{Height functions, the inner product, and the energies}
\label{sec:theory_setup}

Let $\mathcal F=\R^{\text{cells}}$ be the $N$-dimensional space of real
functions on the grid, with the normalised inner product and norm
\begin{equation}
\langle f,g\rangle=\frac1N\sum_{i,j}f_{ij}g_{ij},\qquad \|f\|^{2}=\langle f,f\rangle .
\label{eq:inner}
\end{equation}
The mean of $f$ is $\bar f=\langle f,\mathbf 1\rangle$ and its centring is
$\tilde f=f-\bar f\mathbf 1$.  For two functions (``test function'' $\varphi$ and
height $f$) the empirical covariance over the $N$ cells is
$\Cov(\varphi,f)=\langle\tilde\varphi,\tilde f\rangle=\langle\tilde\varphi,f\rangle=\langle\varphi,\tilde f\rangle$;
the last two equalities hold because $\tilde\varphi\perp\mathbf 1$ and $\tilde f\perp\mathbf 1$.

\begin{lemma}[Means, variances, and covariance as an inner product]
\label{lem:variances}
For every arrangement $S$ of $\{1,\dots,N\}$:
\begin{enumerate}
\item $\bar x=\bar y=\bar z=0$, so $z$ is centred and $\Cov(\varphi,Z)=\langle\tilde\varphi,z\rangle=\langle\varphi,z\rangle$ for every $\varphi\in\mathcal F$;
\item $\Var(X)=\Var(Y)=\dfrac{n^{2}-1}{12}$, \quad $\Var(Z)=\dfrac{N^{2}-1}{12}=\dfrac{n^{4}-1}{12}$, \quad $\Cov(X,Y)=0$;
\item for every line $L$ with indicator $\mathbf 1_L$,\quad $\Cov(\mathbf 1_L,Z)=\dfrac{\delta_L(S)}{N}$.
\end{enumerate}
\end{lemma}

\begin{proof}
(1) $\sum_{i,j}x_{ij}=n\sum_{j}(j-\frac{n-1}2)=0$, likewise for $y$, and $\sum_{i,j}z_{ij}=\sum_{v=1}^{N}v-N\frac{N+1}2=0$;
then $\Cov(\varphi,Z)=\langle\varphi,z\rangle-\bar\varphi\langle\mathbf 1,z\rangle=\langle\varphi,z\rangle$.
(2) $\Var(X)=\frac1n\sum_{j}(j-\frac{n-1}2)^{2}$ is the variance of the uniform distribution on $n$ equally spaced
points, $(n^{2}-1)/12$; the same for $Y$, and for $Z$ with $N$ points.  $\sum_{i,j}x_{ij}y_{ij}=(\sum_jx_{\cdot j})(\sum_iy_{i\cdot})=0$.
(3) By (1), $\Cov(\mathbf 1_L,Z)=\frac1N\sum_{c\in L}z_c=\delta_L/N$ by \eqref{eq:line_dev}.
\end{proof}

We write $R_k=\mathbf 1_{\{i=k\}}$, $C_\ell=\mathbf 1_{\{j=\ell\}}$,
$D_{\mathrm{m}}=\mathbf 1_{\{i=j\}}$, $D_{\mathrm{a}}=\mathbf 1_{\{i+j=n-1\}}$ for the $2n+2$ line indicators.

\begin{definition}[Complete and low-mode energies]
\label{def:energies}
For an arrangement $S$ with height function $z$,
\begin{align}
\Efull(S)&=\sum_{k=0}^{n-1}\Cov(R_k,Z)^{2}+\sum_{\ell=0}^{n-1}\Cov(C_\ell,Z)^{2}
+\Cov(D_{\mathrm m},Z)^{2}+\Cov(D_{\mathrm a},Z)^{2},\label{eq:Efull}\\
\Elow(S)&=\Cov(X,Z)^{2}+\Cov(Y,Z)^{2}+\Cov(D_{\mathrm m},Z)^{2}+\Cov(D_{\mathrm a},Z)^{2}.\label{eq:Elow}
\end{align}
\end{definition}

Note that \eqref{eq:Efull} sums over \emph{all} $n$ rows and \emph{all} $n$ columns.
The word ``energy'' is a label for these sums of squares and nothing more: no Hamiltonian is implied, and the temperature of the annealing in \Cref{sec:algorithms} is a search parameter only.
By \Cref{lem:variances}(3) it has the following elementary description, which we use constantly:
\begin{equation}
N^{2}\Efull(S)=\sum_{L}\delta_L(S)^{2}
=\sum_k\rho_k^{2}+\sum_\ell\gamma_\ell^{2}+\delta_{\mathrm m}^{2}+\delta_{\mathrm a}^{2},
\label{eq:Efull_devs}
\end{equation}
the sum of the squared line-sum deviations over all $2n+2$ lines.  In particular $N^{2}\Efull(S)$ is a
non-negative integer.

\begin{remark}[The four-term energy does not characterise magic squares for $n\ge4$]
\label{rem:four_term_insufficient}
$\Elow=0$ imposes four linear conditions on the centred height functions, while ``magic'' imposes $2n$ independent ones
(\Cref{lem:W_structure}); for $n\ge4$ the two cannot agree even before the permutation constraint is taken into account.
A witness for $n=4$ is
\[
S=\begin{pmatrix}9&4&12&3\\16&1&14&7\\10&15&11&8\\2&5&6&13\end{pmatrix},
\]
whose diagonals both sum to $34$ while the row sums are $(28,38,44,26)$ and the column sums $(37,25,43,31)$; nevertheless $\Cov(X,Z)=\Cov(Y,Z)=0$, so $\Elow(S)=0$ although $\Efull(S)=396/256$.
\Cref{prop:L4} counts all such arrangements, and \Cref{rem:n3_lowmode} explains why none exist for $n=3$.
\end{remark}

The dihedral group $D_4$ acts on the grid by the rotations and reflections of the square, on arrangements by
$(gS)_c=s_{g^{-1}c}$ and on functions by $(gf)_c=f_{g^{-1}c}$, so that $z(gS)=g\,z(S)$ and $\langle gf,gh\rangle=\langle f,h\rangle$.
Every $g$ permutes the $2n$ rows and columns among themselves and the two diagonals among themselves, and acts on the centred
coordinates by a signed permutation: $\{gx,gy\}=\{\pm x,\pm y\}$ as unordered pairs up to sign.

\begin{proposition}[$D_4$-invariance of the energies]
\label{prop:d4_energy}
For every $g\in D_4$ and every arrangement $S$: $\Efull(gS)=\Efull(S)$ and
$\Elow(gS)=\Elow(S)$.
\end{proposition}

\begin{proof}
For any line $L$, $\delta_{gL}(gS)=\sum_{c\in gL}s_{g^{-1}c}=\delta_L(S)$, and $L\mapsto gL$ is a bijection of the set of lines,
so the multiset of deviations is invariant and \eqref{eq:Efull_devs} gives $\Efull(gS)=\Efull(S)$.  For $\Elow$,
$\Cov(X,gZ)=\langle g^{-1}x,z\rangle$ with $g^{-1}x\in\{\pm x,\pm y\}$, and similarly for $y$, so
$\Cov(X,gZ)^{2}+\Cov(Y,gZ)^{2}$ is invariant; the diagonal terms are permuted among themselves.
\end{proof}

An earlier version of this paper omitted the last row and the last column from \eqref{eq:Efull}, on the grounds that
$\sum_k\rho_k=\sum_\ell\gamma_\ell=0$ makes them redundant.  They are redundant for the zero set but not for the values:
the reflection that maps column $n-1$ to column $0$ changes that energy by $(\gamma_{n-1}^{2}-\gamma_{0}^{2})/N^{2}$, so it was not
$D_4$-invariant.  Keeping all $2n+2$ terms (two more than the dimension of the constraint space, \Cref{lem:W_structure}) is the price of
invariance; all $\Efull$ statistics in this paper use \eqref{eq:Efull}.

\subsection{The projection lemma}
\label{sec:theory_projection}

Let $\tilde R_k=R_k-\frac1n\mathbf 1$, $\tilde C_\ell=C_\ell-\frac1n\mathbf 1$,
$\tilde D_{\mathrm m}=D_{\mathrm m}-\frac1n\mathbf 1$, $\tilde D_{\mathrm a}=D_{\mathrm a}-\frac1n\mathbf 1$
be the centred line indicators (each line has $n$ of the $N$ cells, so each mean is $1/n$), and set
\[
W_R=\operatorname{span}\{\tilde R_k\},\quad W_C=\operatorname{span}\{\tilde C_\ell\},\quad
W_D=\operatorname{span}\{\tilde D_{\mathrm m},\tilde D_{\mathrm a}\},\quad
W=W_R+W_C+W_D\subset\mathbf 1^{\perp}.
\]

\begin{lemma}[Structure of $W$]
\label{lem:W_structure}
For $n\ge 3$:
\begin{enumerate}
\item $W_R=\{f:\ f_{ij}=a_i,\ \sum_i a_i=0\}$ and $W_C=\{f:\ f_{ij}=b_j,\ \sum_j b_j=0\}$; both have dimension $n-1$.
\item $\langle\tilde R_k,\tilde R_{k'}\rangle=\frac1N(n\delta_{kk'}-1)$, and the same for columns.
\item $W_R\perp W_C$, $W_R\perp W_D$, $W_C\perp W_D$.
\item $\langle\tilde D_{\mathrm m},\tilde D_{\mathrm m}\rangle=\langle\tilde D_{\mathrm a},\tilde D_{\mathrm a}\rangle=\frac{n-1}{N}$ and
$\langle\tilde D_{\mathrm m},\tilde D_{\mathrm a}\rangle=\frac{\epsilon_n}{N}$ with $\epsilon_n=0$ for $n$ odd and
$\epsilon_n=-1$ for $n$ even; in particular $\tilde D_{\mathrm m},\tilde D_{\mathrm a}$ are linearly independent and $\dim W_D=2$.
\item $W=W_R\oplus W_C\oplus W_D$ is an orthogonal direct sum and $\dim W=2n$.
\end{enumerate}
\end{lemma}

\begin{proof}
For two lines $L,L'$, $N\langle\tilde{\mathbf 1}_L,\tilde{\mathbf 1}_{L'}\rangle=|L\cap L'|-1$, since each line has $n$ of the $N=n^2$ cells.
Two distinct rows are disjoint and a row meets itself in $n$ cells (2); a row meets every column and each diagonal in exactly one
cell, and so does a column (3); the two diagonals meet in the centre cell when $n$ is odd and not at all when $n$ is even (4),
and the Gram matrix $\frac1N\bigl(\begin{smallmatrix}n-1&\epsilon_n\\ \epsilon_n&n-1\end{smallmatrix}\bigr)$ has positive
determinant for $n\ge3$.  For (1), $\sum_ka_k\tilde R_k$ is the function $(i,j)\mapsto a_i-\bar a$, and every zero-mean
function of $i$ alone arises this way; the map $a\mapsto a-\bar a$ has rank $n-1$.  (5) follows: the three subspaces are
mutually orthogonal, hence independent, and $\dim W=2(n-1)+2$.
\end{proof}

For a subspace $V\subset\mathcal F$ let $P_V$ denote the orthogonal projection with respect to \eqref{eq:inner}.

\begin{lemma}[Projection lemma]
\label{lem:projection}
Let $S$ be an arrangement with height function $z$ and line deviations $\rho,\gamma\in\R^{n}$,
$\delta=(\delta_{\mathrm m},\delta_{\mathrm a})\in\R^{2}$, and let $G=\bigl(\begin{smallmatrix}n-1&\epsilon_n\\ \epsilon_n&n-1\end{smallmatrix}\bigr)$
as in \Cref{lem:W_structure}(4).  Then
\begin{enumerate}
\item $P_{W_R}z$ is the function $(i,j)\mapsto\rho_i/n$ (the row means of $z$) and $P_{W_C}z$ is $(i,j)\mapsto\gamma_j/n$, with
$\sum_{k}\Cov(R_k,Z)^{2}=\|\rho\|^{2}/N^{2}=\frac1n\|P_{W_R}z\|^{2}$ and likewise for columns;
\item $\|P_{W_D}z\|^{2}=\frac1N\,\delta^{\mathsf T}G^{-1}\delta$, while $\Cov(D_{\mathrm m},Z)^{2}+\Cov(D_{\mathrm a},Z)^{2}=\delta^{\mathsf T}\delta/N^{2}$;
\item consequently
\begin{equation}
\Efull(S)=\frac1n\Bigl(\|P_{W_R}z\|^{2}+\|P_{W_C}z\|^{2}\Bigr)+\frac{\delta^{\mathsf T}\delta}{N^{2}}=Q\bigl(P_Wz\bigr),
\label{eq:Efull_projection}
\end{equation}
where $Q(w)=\frac1n\|P_{W_R}w\|^{2}+\frac1n\|P_{W_C}w\|^{2}+\frac1{N^{2}}w_D^{\mathsf T}G^{2}w_D$ ($w_D$ the coordinates of
$P_{W_D}w$ in the basis $\tilde D_{\mathrm m},\tilde D_{\mathrm a}$) is a positive-definite quadratic form on $W$.
In particular $\Efull$ depends on $z$ only through $P_Wz$; but $\Efull\ne\|P_Wz\|^{2}$ in general, because the centred
indicators are not an orthonormal system.
\end{enumerate}
\end{lemma}

\begin{proof}
(1) The orthogonal projection of $z$ onto the functions of $i$ alone is the function of row means $(i,j)\mapsto\rho_i/n$, which
has zero mean because $z$ does, so it is $P_{W_R}z$; its squared norm is $\|\rho\|^{2}/n^{3}$ while
$\sum_k\Cov(R_k,Z)^{2}=\|\rho\|^{2}/n^{4}$ by \Cref{lem:variances}(3).  Columns are identical.
(2) The Gram matrix of $(\tilde D_{\mathrm m},\tilde D_{\mathrm a})$ is $G/N$ and the vector of inner products with $z$ is
$\delta/N$, so $\|P_{W_D}z\|^{2}=(\delta/N)^{\mathsf T}(G/N)^{-1}(\delta/N)=\frac1N\delta^{\mathsf T}G^{-1}\delta$.
(3) Sum (1) and (2).  If $P_{W_D}z=w_{\mathrm m}\tilde D_{\mathrm m}+w_{\mathrm a}\tilde D_{\mathrm a}$ then $\delta=Gw_D$,
so $\delta^{\mathsf T}\delta/N^{2}=w_D^{\mathsf T}G^{2}w_D/N^{2}$; each piece of $Q$ is positive definite on its orthogonal
summand ($G$ is invertible), so $Q$ is positive definite on $W$.
\end{proof}

\begin{lemma}[Zero set of the energies]
\label{lem:zero_set}
For every arrangement $S$ of $\{1,\dots,N\}$:
\begin{enumerate}
\item $\Efull(S)=0\iff P_Wz=0\iff z\perp W\iff S$ is a magic square.
\item $\sum_k\Cov(R_k,Z)^{2}+\sum_\ell\Cov(C_\ell,Z)^{2}=0\iff z\perp W_R\oplus W_C\iff S$ is semi-magic.
\item $\Elow(S)=0\implies$ both diagonals sum to $M(n)$; the implication
``$\Elow(S)=0\implies S$ magic'' holds for $n=3$ and fails for $n=4$ (\Cref{rem:n3_lowmode}, \Cref{prop:L4}).
\end{enumerate}
\end{lemma}

\begin{proof}
(1) $\Efull(S)=0$ iff $Q(P_Wz)=0$ iff $P_Wz=0$ iff $z\perp W$, and by \eqref{eq:Efull_devs} iff all line deviations vanish.
(2) is the same with the diagonal terms removed; (3) is immediate from \eqref{eq:Elow} and \Cref{rem:n3_lowmode}.
\end{proof}

\begin{remark}[What the projection lemma does and does not say]
\label{rem:tautology}
\Cref{lem:zero_set}(1) is the definition of a magic square rewritten in the language of orthogonality; it is not a theorem
about magic squares, and we do not present it as one.  Its use is as a change of viewpoint, which makes three things available
that are not visible in the line-sum formulation: (i) $W$ can be replaced by polynomial test functions
(\Cref{sec:theory_hierarchy}), which produces a hierarchy of relaxations whose degree-two level implies Loly's inertia
invariance; (ii) for a uniformly random arrangement the quantities $\langle\varphi,z\rangle$ are permutation statistics in the
sense of Hoeffding, with computable moments and limit law (\Cref{sec:random_energy}); (iii) a single swap changes $z$ by a
multiple of a difference of two cell indicators, and \eqref{eq:Efull_devs} turns the change of $\Efull$ into a two-line count
(\Cref{sec:theory_swap}).
\end{remark}

\subsection{The moment hierarchy}
\label{sec:theory_hierarchy}

The functions $x$ and $y$ are the degree-one polynomials in the column and row index.  Replacing them by
higher powers gives a natural family of relaxations of the semi-magic condition.

\begin{definition}[Degree-$k$ conditions]
\label{def:degree_k}
Let $f\in\mathcal F$ be centred ($\langle f,\mathbf 1\rangle=0$) and $k\ge1$.  We say $f$ satisfies the
\emph{degree-$k$ conditions} if
\[
\langle x^{m},f\rangle=0\quad\text{and}\quad\langle y^{m},f\rangle=0\qquad\text{for all }1\le m\le k .
\]
For $f=z(S)$ these read $\Cov(X^{m},Z)=\Cov(Y^{m},Z)=0$ for $1\le m\le k$, by \Cref{lem:variances}(1).
\end{definition}

Write $x_j=j-\frac{n-1}2$ and $y_i=\frac{n-1}2-i$ for the $n$ distinct values taken by $x$ and $y$, and
for centred $f$ let $\gamma_j(f)=\sum_i f_{ij}$ and $\rho_i(f)=\sum_j f_{ij}$ be its column and row sums
(for $f=z(S)$ these are the deviations $\gamma,\rho$ of \eqref{eq:line_dev}).

\begin{proposition}[Moment hierarchy]
\label{prop:hierarchy}
Let $n\ge3$ and let $f\in\mathcal F$ be centred.
\begin{enumerate}
\item[(a)] For every function $\varphi$ that depends only on the column index, $\langle\varphi,f\rangle=\frac1N\sum_j\varphi(j)\gamma_j(f)$;
in particular $\langle x^{m},f\rangle=\frac1N\sum_j x_j^{m}\gamma_j(f)$ and $\langle y^{m},f\rangle=\frac1N\sum_i y_i^{m}\rho_i(f)$ for all $m\ge0$.
Hence if $\gamma(f)=0$ and $\rho(f)=0$ (for $f=z(S)$: if $S$ is semi-magic), then
$\langle x^{m},f\rangle=\langle y^{m},f\rangle=0$ for \emph{every} $m\ge0$.
\item[(b)] If $k\ge n-1$, the degree-$k$ conditions imply $\gamma(f)=\rho(f)=0$.
\item[(c)] If $1\le k\le n-2$, there is a centred $f$ satisfying the degree-$k$ conditions with $\gamma(f)\ne0$;
explicitly, $f_{ij}=\frac1n(-1)^{j}\binom{k+1}{j}$ for $0\le j\le k+1$ and $f_{ij}=0$ for $j>k+1$.
This $f$ violates the degree-$(k+1)$ condition.
\end{enumerate}
Thus, for real-valued height functions, the degree-$k$ conditions are equivalent to ``equal row sums and equal
column sums'' exactly when $k\ge n-1$.
\end{proposition}

Parts (b) and (c) are the invertibility of the Vandermonde matrix of the nodes $x_0,\dots,x_{n-1}$, and we claim no
novelty for them; the threshold $n-1$ is the benchmark against which \Cref{thm:perm_threshold} is measured.  Part (a) is,
in physical language, the vanishing of the axial multipole moments $\sum_{ij}s_{ij}x_j^{m}$ of a semi-magic charge
distribution, as Rogers and Loly observed for the dipole and quadrupole moments of a magic cube \cite{rogers2005multipole}.  Note that the hierarchy uses
the pure powers $x^{m},y^{m}$ only: the mixed moments $\sum_{ij}s_{ij}x_j^{a}y_i^{b}$ with $a,b\ge1$, beginning with the
product of inertia $\sum_{ij}s_{ij}x_jy_i$, are not determined by the line sums (\Cref{rem:inertia}).

\begin{proof}
(a) $\langle\varphi,f\rangle=\frac1N\sum_j\varphi(j)\sum_i f_{ij}=\frac1N\sum_j\varphi(j)\gamma_j(f)$; the row statement is symmetric.
The last sentence is immediate.

(b) Suppose $\langle x^{m},f\rangle=0$ for $1\le m\le k$ with $k\ge n-1$.  Together with
$\sum_j\gamma_j(f)=\sum_{i,j}f_{ij}=0$ (the case $m=0$), part (a) says that $\gamma(f)\in\R^{n}$ is orthogonal to the
$n$ vectors $(x_j^{m})_{j}$, $m=0,\dots,n-1$.  These are the columns of the Vandermonde matrix of the $n$
\emph{distinct} nodes $x_0,\dots,x_{n-1}$, which is invertible; hence they span $\R^{n}$ and $\gamma(f)=0$.
The same argument with the nodes $y_i$ gives $\rho(f)=0$.

(c) Let $v_j=(-1)^{j}\binom{k+1}{j}$ for $0\le j\le k+1$ and $v_j=0$ for $k+1<j\le n-1$ (this needs
$k+1\le n-1$).  For a polynomial $p$ define the forward difference $(\Delta p)(t)=p(t+1)-p(t)$; $\Delta$
lowers the degree of a nonconstant polynomial by exactly one and kills constants, so $\Delta^{k+1}p=0$ whenever
$\deg p\le k$.  Writing $E$ for the shift $(Ep)(t)=p(t+1)$ we have $\Delta=E-1$, and since $E$ and $1$ commute
the binomial theorem gives $\Delta^{k+1}=\sum_{j=0}^{k+1}\binom{k+1}{j}(-1)^{k+1-j}E^{j}$, i.e.
\begin{equation}
(\Delta^{k+1}p)(0)=(-1)^{k+1}\sum_{j=0}^{k+1}(-1)^{j}\binom{k+1}{j}\,p(j)=(-1)^{k+1}\sum_{j=0}^{n-1}v_j\,p(j).
\label{eq:finite_diff}
\end{equation}
Taking $p(t)=(t-\frac{n-1}2)^{m}$ with $0\le m\le k$ shows $\sum_j v_jx_j^{m}=0$; in particular ($m=0$)
$\sum_jv_j=0$.  Now set $f_{ij}=v_j/n$.  Then $\sum_{i,j}f_{ij}=\sum_jv_j=0$, so $f$ is centred;
$\rho_i(f)=\frac1n\sum_jv_j=0$ for every $i$, so all $\langle y^{m},f\rangle$ vanish by (a);
$\gamma_j(f)=v_j$, so $\langle x^{m},f\rangle=\frac1N\sum_jv_jx_j^{m}=0$ for $1\le m\le k$ by \eqref{eq:finite_diff};
and $\gamma_0(f)=v_0=1\ne0$.  Finally, for $p(t)=(t-\frac{n-1}2)^{k+1}$ we have $\Delta^{k+1}p\equiv(k+1)!\ne0$,
so $\langle x^{k+1},f\rangle=\pm(k+1)!/N\ne0$.
\end{proof}

\begin{corollary}[$\Elow$ is the degree-one truncation]
\label{cor:Elow_degree1}
$\Elow(S)=\langle x,z\rangle^{2}+\langle y,z\rangle^{2}+(\delta_{\mathrm m}^{2}+\delta_{\mathrm a}^{2})/N^{2}$, so
$\Elow(S)=0$ iff $z$ satisfies the degree-$1$ conditions and both diagonals sum to $M(n)$ (\Cref{lem:variances}).
By \Cref{prop:hierarchy}(a) every magic square has $\Elow=0$; by \Cref{prop:hierarchy}(c) the degree-$1$ conditions
characterise equal line sums among real height functions only when $n-1\le1$, i.e.\ never for $n\ge3$.
\end{corollary}

\begin{corollary}[Loly's invariant from the hierarchy]
\label{cor:loly}
Place a mass $s_{ij}$ at the planar point $(x_{ij},y_{ij})$ and let
$I(S)=\sum_{i,j}s_{ij}\bigl(x_{ij}^{2}+y_{ij}^{2}\bigr)$ be the moment of inertia about the centre \cite{loly2004invariance}.
Then for every arrangement
\begin{equation}
I(S)=N\,\Cov\bigl(Z,\,X^{2}+Y^{2}\bigr)+\frac{n^{2}(n^{4}-1)}{12}
=\sum_j x_j^{2}\gamma_j+\sum_i y_i^{2}\rho_i+\frac{n^{2}(n^{4}-1)}{12}.
\label{eq:loly_identity}
\end{equation}
Hence $I(S)$ takes Loly's value $n^{2}(n^{4}-1)/12$ if and only if the single condition $\Cov(Z,X^{2}+Y^{2})=0$ holds.
This condition is implied by the degree-$2$ conditions of \Cref{def:degree_k}, and therefore (\Cref{prop:hierarchy}(a)) by the
semi-magic property, but it is much weaker than the degree-$2$ conditions (see below).
The two halves $I_x=\sum_jx_j^{2}(\gamma_j+M(n))$ and $I_y=\sum_iy_i^{2}(\rho_i+M(n))$ each equal $n^{2}(n^{4}-1)/24$ for
semi-magic squares.
\end{corollary}

The invariance and its validity for semi-magic squares are Loly's \cite{loly2004invariance}: his derivation writes
$I=I_x+I_y$ with $I_x=\sum_jx_j^{2}\cdot(\text{column sum }j)$, $I_y=\sum_iy_i^{2}\cdot(\text{row sum }i)$, uses the row and
column sums only, and gives $I_x=I_y=n^{2}(n^{4}-1)/24$.  The corollary re-derives this; what it adds is the decomposition
\eqref{eq:loly_identity} into a degree-two cross-moment plus a constant, and the observation that the invariance is one scalar
condition rather than the degree-two level.  How much weaker the single condition is can be seen for $n=3$ by enumeration
of all $9!$ arrangements (\texttt{check\_round1.py}): $24{,}192$ of them have $I=60$, against $72$ satisfying the
degree-$2$ conditions (the semi-magic squares, \Cref{rem:n3_lowmode}).  The values for $n=3,\dots,9$ are $60$, $340$, $1300$, $3885$, $9800$, $21840$,
$44280$ (OEIS~A126275).

\begin{proof}
Substitute $s_{ij}=z_{ij}+\frac{N+1}2$:
$I(S)=\sum_{i,j}z_{ij}(x_{ij}^{2}+y_{ij}^{2})+\frac{N+1}2\sum_{i,j}(x_{ij}^{2}+y_{ij}^{2})$.
The first sum is $N\Cov(X^{2}+Y^{2},Z)$ by \Cref{lem:variances}(1) and equals
$\sum_jx_j^{2}\gamma_j+\sum_iy_i^{2}\rho_i$ by \Cref{prop:hierarchy}(a); the second is
$\frac{N+1}2\cdot2\cdot\frac{n^{2}(n^{2}-1)}{12}=\frac{n^{2}(n^{4}-1)}{12}$ by \Cref{lem:variances}(2).
For semi-magic squares $\gamma=\rho=0$ and $I_x=M(n)\sum_jx_j^{2}=\frac{n(n^{2}+1)}2\cdot\frac{n(n^{2}-1)}{12}$.
\end{proof}

\begin{corollary}[Null distribution of Loly's invariant]
\label{cor:loly_null}
For a uniformly random arrangement $S$ of order $n$,
\begin{equation}
\E[I(S)]=\frac{n^{2}(n^{4}-1)}{12},\qquad
\Var\bigl(I(S)\bigr)=N^{2}\sigma_n^{2}\,\Var(X^{2}+Y^{2})=\frac{n^{4}(n^{4}-1)(n^{2}-4)}{1080},
\label{eq:loly_null}
\end{equation}
where $\sigma_n^{2}=(n^{2}+1)/12$ and $\Var(X^{2}+Y^{2})=(n^{2}-1)(n^{2}-4)/90$ is the variance of $x^{2}+y^{2}$ over the $N$ cells.
The mean equals Loly's value for every distribution of arrangements that is invariant under relabelling of the values
$1,\dots,N$ (in particular the uniform distribution):
magic and semi-magic squares sit exactly at the population mean of $I$, and the invariance says that their fluctuation
vanishes, not that their value is extreme.  For $n=3$, $\Var(I)=30$ (standard deviation $5.48$; $29$ values from $46$ to $74$);
for $n=4$, $\Var(I)=2176/3$ (standard deviation $26.9$), and $2.9\%$ of arrangements have $I=340$.
\end{corollary}

\begin{proof}
$N\Cov(Z,X^{2}+Y^{2})=\sum_c\bigl(x_c^{2}+y_c^{2}-\overline{x^{2}+y^{2}}\bigr)z_{\pi(c)}$ is a linear permutation statistic
with mean $0$ (\Cref{sec:bg_perm}); by \eqref{eq:ww_identity} with $a=x^{2}+y^{2}$ and $b=z$ its variance is
$\frac1{N-1}\cdot N\Var(X^{2}+Y^{2})\cdot N\Var(Z)=N^{2}\sigma_n^{2}\Var(X^{2}+Y^{2})$, as in \Cref{prop:random_energy_moments}.
Since $X$ and $Y$ are independent and uniform on $n$ equally spaced points, $\Var(X^{2}+Y^{2})=2\Var(X^{2})=2(\E X^{4}-(\E X^{2})^{2})$
with $\E X^{2}=(n^{2}-1)/12$ and $\E X^{4}=(n^{2}-1)(3n^{2}-7)/240$, which gives $(n^{2}-1)(n^{2}-4)/90$.
The $n=3$ values are confirmed by exhaustive enumeration and the $n=4$ values by $2\times10^{7}$ Monte Carlo samples.
\end{proof}

\Cref{prop:hierarchy} is a statement about real-valued height functions.  For arrangements, i.e.\ for
$f=z(S)$ with $S$ a permutation of $\{1,\dots,N\}$, the integrality and boundedness of the line sums
change the picture in an interesting way.

\begin{theorem}[Integrality lowers the permutation threshold]
\label{thm:perm_threshold}
For $n\ge3$ let $k^{*}(n)$ be the least $k$ such that every arrangement $S$ of $\{1,\dots,N\}$ satisfying
the degree-$k$ conditions is semi-magic.  Then $k^{*}(n)\le n-2$ for all $n\ge14$: for such $n$ the degree-$(n-2)$
conditions already force equal row and column sums.
By contrast $k^{*}(n)\le n-1$ for all $n$ (\Cref{prop:hierarchy}(b)), with equality for $n=3,4$: $k^{*}(3)=2$ and $k^{*}(4)=3$
by the witnesses
\begin{align*}
n=3,\ \text{degree }1:&\quad
\begin{pmatrix}1&3&9\\8&6&5\\7&4&2\end{pmatrix},\qquad \gamma=(1,-2,1),\ \rho=(-2,4,-2);\\[2pt]
n=4,\ \text{degree }2:&\quad
\begin{pmatrix}14&2&5&13\\8&15&7&4\\12&3&9&10\\1&11&16&6\end{pmatrix},\qquad \gamma=(1,-3,3,-1),\ \rho=0.
\end{align*}
\end{theorem}

\begin{proof}
\emph{Lower bounds.}  For the $3\times3$ array, the column sums are $16,13,16$ and the row sums $13,19,13$, so
$\gamma=(1,-2,1)$, $\rho=(-2,4,-2)$ and $\sum_jx_j\gamma_j=(-1)(1)+0+(1)(1)=0$, $\sum_iy_i\rho_i=(1)(-2)+0+(-1)(-2)=0$:
the degree-$1$ conditions hold but the square is not semi-magic, so $k^{*}(3)\ge2$.
For the $4\times4$ array, every row sums to $34$ and the column sums are $35,31,37,33$, so $\rho=0$ and
$\gamma=(1,-3,3,-1)$, which is the vector $v$ of \Cref{prop:hierarchy}(c) with $k=2$; hence
$\sum_jx_j^{m}\gamma_j=0$ for $m=1,2$ by \eqref{eq:finite_diff}, the degree-$2$ conditions hold, and the square
is not semi-magic, so $k^{*}(4)\ge3$.

\emph{The bound for $n\ge14$.}  Let $S$ satisfy the degree-$(n-2)$ conditions.  By \Cref{prop:hierarchy}(a), $\gamma$ is orthogonal to the
$n-1$ linearly independent vectors $(x_j^{m})_j$, $m=0,\dots,n-2$, so it lies in their one-dimensional orthogonal
complement in $\R^{n}$.  By \eqref{eq:finite_diff} with $k+1=n-1$, this complement is spanned by
$v_j=(-1)^{j}\binom{n-1}{j}$, $j=0,\dots,n-1$ (which is orthogonal to all polynomials of degree $\le n-2$ and
nonzero).  So $\gamma=tv$ with $t=\gamma_0\in\Z$ because $v_0=1$ and $\gamma$ is an integer vector.
On the other hand every column sum of a permutation of $\{1,\dots,N\}$ lies between the sum of the $n$ smallest
and the sum of the $n$ largest values, so
\[
|\gamma_j|\le\sum_{v=N-n+1}^{N}v-M(n)=\frac{n(2N-n+1)}{2}-\frac{n(N+1)}{2}=\frac{n(N-n)}{2}=\frac{n^{2}(n-1)}{2}.
\]
If $t\ne0$ then $\max_j|\gamma_j|\ge\binom{n-1}{\lfloor(n-1)/2\rfloor}$.  We claim
$\binom{n-1}{\lfloor(n-1)/2\rfloor}>\frac{n^{2}(n-1)}2$ for all $n\ge14$.  For $n=14$: $\binom{13}{6}=1716>1274=\frac{14^{2}\cdot13}{2}$.
Passing from $n$ to $n+1$ multiplies the left side by $\binom{n}{\lfloor n/2\rfloor}/\binom{n-1}{\lfloor(n-1)/2\rfloor}$,
which equals $2$ for $n$ even and $\frac{2n}{n+1}\ge\frac{28}{15}$ for odd $n\ge14$, while it multiplies the right side by
$\frac{(n+1)^{2}n}{n^{2}(n-1)}=\frac{(n+1)^{2}}{n(n-1)}\le\frac{225}{182}<\frac{28}{15}$ for $n\ge14$; the claim follows by induction.
Hence $t=0$ and $\gamma=0$; the same argument applied to the rows gives $\rho=0$.
\end{proof}

\begin{remark}[$k^{*}(n)=n-1$ for $5\le n\le13$; computer-verified]
\label{rem:kstar_search}
For each $5\le n\le13$ a randomised swap search found an arrangement with all row sums equal to $M(n)$ and column deviations
$\gamma_j=(-1)^{j}\binom{n-1}{j}$; by \eqref{eq:finite_diff} such an arrangement satisfies the degree-$(n-2)$ conditions
without being semi-magic, so $k^{*}(n)=n-1$ for all $3\le n\le13$.  The $n=5$ witness is printed in \Cref{sec:app_witnesses};
all nine are in the ancillary file \texttt{revision\_round2/appendix\_data.json}, verified in exact integer arithmetic by
\texttt{check\_round1.py}.  Thus the permutation threshold agrees with the linear-algebra threshold exactly for $n\le13$ and drops
below it from $n=14$ on; determining $k^{*}(n)$ for $n\ge14$ is open (\Cref{sec:discussion}).
\end{remark}

\begin{remark}[Why the low-mode energy appeared to work for $n=3$]
\label{rem:n3_lowmode}
Exhaustive enumeration of the $9!$ arrangements (\texttt{verify\_theory.py}, exact arithmetic) gives $760$ arrangements
satisfying the degree-$1$ conditions, exactly $72$ (the semi-magic squares) satisfying the degree-$2$ conditions, in accordance
with $k^{*}(3)=2$, and exactly $8$ satisfying the degree-$1$ conditions together with both diagonal conditions---the magic
squares.  So for $n=3$, $\Elow(S)=0\iff S$ is magic, but for the ``wrong'' reason: the degree-$1$ conditions leave the
families $\gamma=t(1,-2,1)$, $\rho=u(1,-2,1)$ open, and the two diagonal conditions combined with the permutation constraint
happen to force $t=u=0$.  Nothing of the sort survives for $n=4$: the arrangement of \Cref{rem:four_term_insufficient} has
$\Elow=0$ with $\gamma=(3,-9,9,-3)$, $\rho=(-6,4,10,-8)$; it satisfies the degree-$2$ condition in $x$ but not in $y$
($\sum_iy_i^{2}\rho_i=-28$).
\end{remark}

\subsection{The single-swap formula}
\label{sec:theory_swap}

Let $a\ne b$ be two cells and let $S^{ab}$ be the arrangement obtained from $S$ by exchanging the entries in
$a$ and $b$.  Define
\begin{align*}
m(a,b)&=\#\{\text{lines }L:\ L\text{ contains exactly one of }a,b\},\\
u_S(c)&=\sum_{L\ni c}\delta_L(S)\qquad(\text{sum over the lines through }c).
\end{align*}

\begin{theorem}[Single-swap formula]
\label{thm:swap}
For every arrangement $S$ and cells $a\ne b$, with $d=s_b-s_a$,
\begin{equation}
N^{2}\bigl(\Efull(S^{ab})-\Efull(S)\bigr)=d^{2}\,m(a,b)+2d\,\bigl(u_S(a)-u_S(b)\bigr).
\label{eq:swap_general}
\end{equation}
In particular, if $S$ is a magic square then, with $d=|s_a-s_b|$,
\begin{equation}
\Efull(S^{ab})=\frac{d^{2}\,m(a,b)}{N^{2}} .
\label{eq:swap_magic}
\end{equation}
\end{theorem}

\begin{proof}
The swap changes the height function by $z(S^{ab})=z(S)+d(\mathbf 1_a-\mathbf 1_b)$.  For a line $L$,
$\delta_L(S^{ab})=\delta_L(S)+d\bigl(\mathbf 1_L(a)-\mathbf 1_L(b)\bigr)$: the deviation increases by $d$ if $L$
contains $a$ but not $b$, decreases by $d$ if $L$ contains $b$ but not $a$, and is unchanged if $L$ contains both
or neither.  By \eqref{eq:Efull_devs},
\begin{align*}
N^{2}\Delta \Efull&=\sum_{L\ni a,\,L\not\ni b}\bigl[(\delta_L+d)^{2}-\delta_L^{2}\bigr]
+\sum_{L\ni b,\,L\not\ni a}\bigl[(\delta_L-d)^{2}-\delta_L^{2}\bigr]\\
&=d^{2}m(a,b)+2d\Bigl(\sum_{L\ni a,\,L\not\ni b}\delta_L-\sum_{L\ni b,\,L\not\ni a}\delta_L\Bigr).
\end{align*}
Adding and subtracting the lines containing both $a$ and $b$ turns the bracket into $u_S(a)-u_S(b)$.
If $S$ is magic, all $\delta_L(S)=0$, so $u_S\equiv0$ and $\Efull(S)=0$, giving \eqref{eq:swap_magic}.
\end{proof}

\begin{lemma}[The multiplicity $m(a,b)$]
\label{lem:m_values}
For distinct cells $a=(i,j)$, $b=(i',j')$:
\begin{align*}
m(a,b)=2\cdot[i\ne i']+2\cdot[j\ne j']&+[\text{exactly one of }a,b\text{ on the main diagonal}]\\
&+[\text{exactly one of }a,b\text{ on the anti-diagonal}].
\end{align*}
Hence $2\le m(a,b)\le6$, and $m(a,b)=2$ if and only if $a$ and $b$ lie in a common row or a common column and
neither of them lies on a diagonal.  Such pairs exist for every $n\ge3$ (for $n=3$ exactly one pair up to $D_4$:
the two non-central cells of the middle row or column; for $n\ge4$ any two off-diagonal cells of a row).
\end{lemma}

\begin{proof}
If $i=i'$ the common row contains both cells and no other row contains either; if $i\ne i'$ the row of $a$ contains $a$ only
and the row of $b$ contains $b$ only; columns likewise, and each diagonal contributes $1$ or $0$ according to whether it
separates $a$ from $b$.  Since $a\ne b$ the row/column contribution is $2$ or $4$, so $2\le m\le 6$.  If $m=2$ then $i=i'$ or
$j=j'$ (not both) and no diagonal separates the cells; for $i=i'$, $j\ne j'$ both cells cannot lie on the same diagonal, so
``not separated'' means neither is on it.  For existence: a row contains at most two diagonal cells, so it has at least
$n-2\ge2$ off-diagonal cells when $n\ge4$, and the middle row of the $3\times3$ grid has the two off-diagonal cells $(1,0),(1,2)$.
\end{proof}

\begin{corollary}[Strict local minimality and the exact gap]
\label{cor:local_min}
Let $S$ be a magic square of order $n\ge3$.
\begin{enumerate}
\item Every single swap strictly increases the energy: $\Efull(S^{ab})\ge 2/N^{2}>0$ for all $a\ne b$.
Thus every magic square is a strict local minimum of $\Efull$ on the set of arrangements with respect to transpositions, for every $n$.
\item The smallest energy reachable by one swap is
\[
\Delta_{\min}(S)=\min_{a\ne b}\frac{(s_a-s_b)^{2}\,m(a,b)}{N^{2}},\qquad
\frac{2}{N^{2}}\le\Delta_{\min}(S)\le\frac{6}{N^{2}} .
\]
The lower bound is attained iff two cells with consecutive values share a row or column and are both off the diagonals;
the upper bound is the swap of any two consecutive values ($d=1$, $m\le6$).  Thus $\Delta_{\min}(S)=\Theta(n^{-4})$ with
explicit constants, uniformly over all magic squares of order $n$.
\item The possible values of $N^{2}\Efull(S^{ab})$ are the integers $d^{2}m$ with $1\le d\le N-1$ and $2\le m\le6$.
For the $880$ magic squares of order $4$ the multiset of $N^{2}\Delta_{\min}$ is $\{2:432,\ 3:420,\ 4:28\}$, and for the
$8$ squares of order $3$ it is $\{3:8\}$ (\Cref{sec:swap-gaps}).
\end{enumerate}
\end{corollary}

\begin{proof}
By \eqref{eq:swap_magic}, $N^{2}\Efull(S^{ab})=d^{2}m(a,b)$ with $d\ge1$ and $m\ge2$ (\Cref{lem:m_values}); the cells
carrying the values $v,v+1$ have $d=1$ and $m\le6$, and the equality case of the lower bound is \Cref{lem:m_values}.
\end{proof}

\begin{remark}[Swap formula for $\Elow$]
\label{rem:swap_low}
Since $\Cov(\varphi,Z^{ab})=\Cov(\varphi,Z)+\frac dN(\varphi(a)-\varphi(b))$ for every test function $\varphi$
(\Cref{lem:variances}(1)), the same computation gives, for $\Elow(S)=0$ and in particular for $S$ magic,
\[
\Elow(S^{ab})=\frac{d^{2}}{N^{2}}\Bigl[(x_a-x_b)^{2}+(y_a-y_b)^{2}+\bigl(D_{\mathrm m}(a)-D_{\mathrm m}(b)\bigr)^{2}+\bigl(D_{\mathrm a}(a)-D_{\mathrm a}(b)\bigr)^{2}\Bigr],
\]
i.e.\ $d^{2}/N^{2}$ times the squared planar distance between the cells plus the number of diagonals separating them.
The bracket is at least $1$, so magic squares are strict local minima of $\Elow$ as well, with $\Elow(S^{ab})\ge1/N^{2}$.
\end{remark}

\begin{remark}[The inertia tensor of the gem]
\label{rem:inertia}
Place unit masses at the $N$ points $\mathbf p_{ij}$.  The centre of mass is the origin (\Cref{lem:variances}(1)) and, by
\eqref{eq:inertia_cov} and \Cref{lem:variances}(2), the inertia tensor has
$I_{xx}=I_{yy}=N(\Var Y+\Var Z)=n^{2}(n^{2}-1)(n^{2}+2)/12$, $I_{zz}=N(\Var X+\Var Y)=n^{2}(n^{2}-1)/6$, $I_{xy}=-N\Cov(X,Y)=0$,
$I_{xz}=-N\Cov(X,Z)=-\sum_jx_j\gamma_j$ and $I_{yz}=-\sum_iy_i\rho_i$ (\Cref{prop:hierarchy}(a)).
The diagonal entries are the same for all $N!$ arrangements, and for every arrangement satisfying the degree-$1$ conditions the
tensor is diagonal: the unit-mass gem is an oblate symmetric top whose inertia carries no information about the square.
The informative planar object is Loly's value-weighted one: its diagonal entries $\sum s_{ij}y_{ij}^{2}$, $\sum s_{ij}x_{ij}^{2}$
are fixed by the row and column sums (\Cref{cor:loly}), whereas the product of inertia $P_{xy}=-\sum_{ij}s_{ij}x_{ij}y_{ij}$ is a
mixed moment, fixed by no line sum: $D_4$ changes at most its sign, and over the $880$ essentially different squares of order $4$
the invariant $|P_{xy}|$ takes $48$ distinct values between $0$ and $60$, vanishing for $56$ of them.
\end{remark}
